\documentclass[a4paper,10pt,BCOR10mm,oneside,headsepline]{scrartcl}
\usepackage[portuges]{babel}
\usepackage[utf8]{inputenc}
\usepackage{amsmath}
\usepackage{wasysym}% provides \ocircle and \Box
\usepackage{enumitem}% easy control of topsep and leftmargin for lists
\usepackage{color}% used for background color
\usepackage{forloop}% used for \Qrating and \Qlines
\usepackage{ifthen}% used for \Qitem and \QItem
\usepackage{typearea}
\usepackage{graphicx}
\usepackage{booktabs}
\usepackage{listings}
\usepackage{color}
\usepackage{algorithm}
\usepackage{algorithmic}
\usepackage{amssymb}
\usepackage{minted}
\usepackage{mdframed}
\usepackage{wasysym}
\usepackage{xcolor}

\usepackage{tgtermes}

\areaset{17cm}{26cm}
\setlength{\topmargin}{-2cm}

\definecolor{codegreen}{rgb}{0,0.6,0}
\definecolor{codegray}{rgb}{0.5,0.5,0.5}
\definecolor{codepurple}{rgb}{0.58,0,0.82}
\definecolor{backcolour}{rgb}{0.95,0.95,0.92}
 
\lstdefinestyle{mystyle}{
    backgroundcolor=\color{backcolour},   
    commentstyle=\color{codegreen},
    keywordstyle=\color{magenta},
    numberstyle=\tiny\color{codegray},
    stringstyle=\color{codepurple},
    basicstyle=\footnotesize,
    breakatwhitespace=false,         
    breaklines=true,                 
    captionpos=b,                    
    keepspaces=true,                 
    numbers=left,                    
    numbersep=5pt,                  
    showspaces=false,                
    showstringspaces=false,
    showtabs=false,                  
    tabsize=2
}
 
\lstset{style=mystyle}

\mdfsetup{
    skipabove=12pt,
    skipbelow=12pt,
    innertopmargin=8pt,
    innerbottommargin=8pt,
    innerleftmargin=10pt,
    innerrightmargin=10pt,
    splittopskip=12pt,
    splitbottomskip=12pt,
}

\algsetup{linenodelimiter=.,linenosize=\tiny}

\renewcommand{\algorithmicrequire}{\textbf{Entrada:}}
\renewcommand{\algorithmicensure}{\textbf{Saída:}}
\floatname{algorithm}{Algoritmo}

%% \Qq = Questionaire question. Oh, this is just too simple. It helps
%% making it easy to globally change the appearance of questions.
\newcommand{\Qq}[1]{\textbf{#1}}
%\newcommand{\Qq}[1]{{#1}}

%% \QO = Circle or box to be ticked. Used both by direct call and by
%% \Qrating and \Qlist.
\newcommand{\QO}{$\Box$}% or: $\ocircle$

%% \Qrating = Automatically create a rating scale with NUM steps, like
%% this: 0--0--0--0--0.
\newcounter{qr}
\newcommand{\Qrating}[1]{\QO\forloop{qr}{1}{\value{qr} < #1}{---\QO}}

%% \Qline = Again, this is very simple. It helps setting the line
%% thickness globally. Used both by direct call and by \Qlines.
\newcommand{\Qline}[1]{\noindent\rule{#1}{0.6pt}}

%% \Qlines = Insert NUM lines with width=\linewith. You can change the
%% \vskip value to adjust the spacing.
\newcounter{ql}
\newcommand{\Qlines}[1]{\forloop{ql}{0}{\value{ql}<#1}{\vskip0em\Qline{\linewidth}}}

%% \Qlist = This is an environment very similar to itemize but with
%% \QO in front of each list item. Useful for classical multiple
%% choice. Change leftmargin and topsep accourding to your taste.
\newenvironment{Qlist}{%
\renewcommand{\labelitemi}{\QO}
\begin{itemize}[leftmargin=1.5em,topsep=-.5em]
}{%
\end{itemize}
}

%% \Qtab = A "tabulator simulation". The first argument is the
%% distance from the left margin. The second argument is content which
%% is indented within the current row.
\newlength{\qt}
\newcommand{\Qtab}[2]{
\setlength{\qt}{\linewidth}
\addtolength{\qt}{-#1}
\hfill\parbox[t]{\qt}{\raggedright #2}
}

%% \Qitem = Item with automatic numbering. The first optional argument
%% can be used to create sub-items like 2a, 2b, 2c, ... The item
%% number is increased if the first argument is omitted or equals 'a'.
%% You will have to adjust this if you prefer a different numbering
%% scheme. Adjust topsep and leftmargin as needed.
\newcounter{itemnummer}
\newcommand{\Qitem}[2][]{% #1 optional, #2 notwendig
\ifthenelse{\equal{#1}{}}{\stepcounter{itemnummer}}{}
\ifthenelse{\equal{#1}{a}}{\stepcounter{itemnummer}}{}
\begin{enumerate}[topsep=2pt,leftmargin=2.8em]
\item[\textbf{\arabic{itemnummer}#1.}] #2
\end{enumerate}
}

%% \QItem = Like \Qitem but with alternating background color. This
%% might be error prone as I hard-coded some lengths (-5.25pt and
%% -3pt)! I do not yet understand why I need them.
\definecolor{bgodd}{rgb}{0.8,0.8,0.8}
\definecolor{bgeven}{rgb}{0.9,0.9,0.9}
\newcounter{itemoddeven}
\newlength{\gb}
\newcommand{\QItem}[2][]{% #1 optional, #2 notwendig
\setlength{\gb}{\linewidth}
\addtolength{\gb}{-5.25pt}
\ifthenelse{\equal{\value{itemoddeven}}{0}}{%
\noindent\colorbox{bgeven}{\hskip-3pt\begin{minipage}{\gb}\Qitem[#1]{#2}\end{minipage}}%
\stepcounter{itemoddeven}%
}{%
\noindent\colorbox{bgodd}{\hskip-3pt\begin{minipage}{\gb}\Qitem[#1]{#2}\end{minipage}}%
\setcounter{itemoddeven}{0}%
}
}

\newcommand{\rawcode}[1]{\texttt{#1}}
\newcommand{\code}[2]{\mintinline{#1}{#2}}
\newcommand{\file}[1]{\texttt{\underline{#1}}}

\renewcommand{\theFancyVerbLine}{\tiny \texttt{\textcolor{gray}{\arabic{FancyVerbLine}}}}

\newcommand{\inputcode}[2]{\begin{mdframed}\inputminted[fontsize=\footnotesize,linenos,numbersep=4pt]{#1}{#2}\end{mdframed}}

\newenvironment{solution}{%
    \begin{mdframed}%
        \noindent\textbf{Solu\c{c}\~ao:}\par\medskip
}{%
    \end{mdframed}%
}

\newcommand{\inputsnippet}[4]{\begin{mdframed}[shadow=true,shadowsize=3pt,shadowcolor=gray]\inputminted[fontsize=\footnotesize,linenos,firstline=#2,lastline=#3,numbersep=2pt]{#1}{#4}\end{mdframed}}

\newcommand{\inputsyntax}[2]{\inputminted[fontsize=\footnotesize]{#1}{#2}}



\begin{document}

\begin{center}
	\includegraphics[scale=0.45]{fga}
	% fga.ps: 1087x103 pixel, 72dpi, 38.35x3.63 cm, bb=0 0 1087 103
\end{center}


\vspace{0.1in}

\begin{center}
    \textbf{\huge Paradigmas de Programação}

    \vspace{0.1in}
    \textbf{\Large Cálculo-$\lambda$: Lista 2}
\end{center}\vskip1em

\renewcommand{\Qq}[1]{{#1}}

\section*{Parte A - Aritmética}

\Qitem{ \Qq{Defina um termo-$\lambda$ $L$ tal que, para dois naturais $x$ e $y$, $Lxy$ seja
        verdadeiro se $x\leq y$, e falso, caso contrário. Use, nesta definição, os termos-$\lambda$
        $Z$ e $P$.} }

\begin{solution}
Lembre que um natural $n$ aplica $n$ vezes seu primeiro argumento ao segundo. Assim,
$yPx$ aplica $y$ vezes o antecessor a $x$, resultando em $x - y$ se $x\geq y$, ou em $0$
caso contrário (pois $P0=0$). Logo, $x\leq y$ se, e somente se, $yPx = 0$:
\[
    L \;\equiv\; \lambda xy.\,Z(yPx).
\]

\textbf{Verificação:}
\begin{align*}
    L\,2\,3 &\equiv Z(3P2) \equiv Z(P(P(P\,2))) \equiv Z\,0 \equiv T, \\
    L\,3\,3 &\equiv Z(3P3) \equiv Z\,0 \equiv T, \\
    L\,3\,2 &\equiv Z(2P3) \equiv Z(P(P\,3)) \equiv Z\,1 \equiv F.
\end{align*}
(Note que nos slides, há o termo $G$, o termo $L$ é ele com os argumentos trocados: $Lxy \equiv Gyx$.)
\end{solution}

\Qitem{ \Qq{Defina um termo-$\lambda$ $M$ tal que, para dois naturais $x$ e $y$, $Mxy$ seja
        igual ao maior dentre ambos. Use, nesta definição, os termos-$\lambda$ $Z$ e $P$. } }

\begin{solution}
Usando a estrutura \texttt{if-then-else} ($I_F\,p\,a\,b\equiv p\,a\,b$) com o predicado
$x\leq y$, isto é, $Z(yPx)$ da questão anterior: se $x\leq y$, o maior é $y$; caso
contrário, é $x$.
\[
    M \;\equiv\; \lambda xy.\,Z(yPx)\,y\,x.
\]

\textbf{Verificação:}
\begin{align*}
    M\,2\,3 &\equiv Z(3P2)\,3\,2 \equiv Z\,0\,3\,2 \equiv T\,3\,2 \equiv 3, \\
    M\,3\,2 &\equiv Z(2P3)\,2\,3 \equiv Z\,1\,2\,3 \equiv F\,2\,3 \equiv 3.
\end{align*}
\end{solution}

\Qitem{ \Qq{Defina um termo-$\lambda$ $N$ tal que, para dois naturais $x$ e $y$, $N xy$ seja
        igual ao menor dentre ambos. Use, nesta definição, os termos-$\lambda$ $Z$ e $P$. } }

\begin{solution}
Basta trocar a posição do resultado do \texttt{if-then-else} do termo anterior: se $x\leq y$, o menor é $x$; caso contrário, é $y$.
\[
    N \;\equiv\; \lambda xy.\,Z(yPx)\,x\,y.
\]

\textbf{Verificação:}
\begin{align*}
    N\,2\,3 &\equiv Z(3P2)\,2\,3 \equiv Z\,0\,2\,3 \equiv T\,2\,3 \equiv 2, \\
    N\,3\,2 &\equiv Z(2P3)\,3\,2 \equiv Z\,1\,3\,2 \equiv F\,3\,2 \equiv 2.
\end{align*}
\end{solution}

\section*{Parte B - Recursão}

\Qitem{ \Qq{No cálculo lambda, a adição de dois números naturais $x$ e $y$ pode ser definida por meio
        do termo-$\lambda$
    \[
        + \equiv \lambda xy.xSy,
    \]
    onde $Sn$ é o sucessor do natural $n$. Também é possível definir a adição nos naturais por
    meio de recursão:
    \[
        \left\{ \begin{array}{ll} x + 1 = x'\\ x + y' = (x + y)'\end{array}\right.
    \]
    onde $n'$ é o sucessor do natural $n$.

    Usando os termos-$\lambda$ sucessor $S$, antecessor $P$ e igualdade $E$, defina um termo 
    $\Gamma$ tal que $Y\Gamma$ corresponda à esta definição recursiva de adição de naturais. } }

\begin{solution}
Seguindo a estrutura $\Gamma\equiv\lambda f x.(\text{predicado})(\text{caso base})(\text{chamada recursiva})$
nos slides, a recursão é feita sobre $y$:
\begin{itemize}
    \item predicado do caso base: $y = 1$, isto é, $E\,y\,1$;
    \item caso base: $x + 1 = x'$, isto é, $S\,x$;
    \item chamada recursiva: $x + y = (x + (y-1))'$, isto é, $S\,(f\,x\,(P\,y))$.
\end{itemize}
\[
    \Gamma \;\equiv\; \lambda fxy.\,(E\,y\,1)\,(S\,x)\,\bigl(S\,(f\,x\,(P\,y))\bigr).
\]

Considerando que o menor natural é $1$.

\textbf{Verificação} ($2 + 2$):
\begin{align*}
    (Y\Gamma)\,2\,2 &\equiv \Gamma\,(Y\Gamma)\,2\,2
        \equiv (E\,2\,1)\,(S\,2)\,\bigl(S\,((Y\Gamma)\,2\,(P\,2))\bigr)
        \equiv F\,(S\,2)\,\bigl(S\,((Y\Gamma)\,2\,1)\bigr) \\
    &\equiv S\,\bigl((Y\Gamma)\,2\,1\bigr)
        \equiv S\,\bigl(\Gamma\,(Y\Gamma)\,2\,1\bigr)
        \equiv S\,\bigl((E\,1\,1)\,(S\,2)\,(\ldots)\bigr)
        \equiv S\,(T\,(S\,2)\,(\ldots)) \\
    &\equiv S\,(S\,2) \equiv S\,3 \equiv 4.
\end{align*}

\end{solution}

\Qitem{ \Qq{No cálculo lambda, a multiplicação de dois números naturais $x$ e $y$ pode ser definida por
        meio do termo-$\lambda$
    \[
        \times \equiv \lambda xyz.x(yz)
    \]
    Também é possível definir a multiplicação nos naturais por meio de recursão:
    \[
        \left\{ \begin{array}{ll} x\times 1 = x\\ x\times y' = x\times y + x\end{array}\right.
    \]
    onde $n'$ é o sucessor do natural $n$.

    Usando os termos-$\lambda$ sucessor $S$, antecessor $P$, igualdade $E$ e soma $+$, defina um 
    termo $\Gamma$ tal que $Y\Gamma$ corresponda à esta definição recursiva de multiplicação de
    naturais. } }

\begin{solution}
A recursão é feita sobre $y$:
\begin{itemize}
    \item predicado do caso base: $E\,y\,1$;
    \item caso base: $x\times 1 = x$, isto é, $x$;
    \item chamada recursiva: $x\times y = x\times(y-1) + x$, isto é, $+\,(f\,x\,(P\,y))\,x$.
\end{itemize}
\[
    \Gamma \;\equiv\; \lambda fxy.\,(E\,y\,1)\,x\,\bigl(+\,(f\,x\,(P\,y))\,x\bigr).
\]

Considerando que o menor natural é $1$.

\textbf{Verificação} ($3\times 2$):
\begin{align*}
    (Y\Gamma)\,3\,2 &\equiv \Gamma\,(Y\Gamma)\,3\,2
        \equiv (E\,2\,1)\,3\,\bigl(+\,((Y\Gamma)\,3\,(P\,2))\,3\bigr)
        \equiv +\,((Y\Gamma)\,3\,1)\,3 \\
    &\equiv +\,\bigl((E\,1\,1)\,3\,(\ldots)\bigr)\,3
        \equiv +\,3\,3 \equiv 6.
\end{align*}

\end{solution}

\Qitem{ \Qq { Defina um termo-$\lambda$ $\Gamma$ tal que $(Y\Gamma)$ compute, recursivamente, o que
        for solicitado nos itens a seguir. É permitido o uso de qualquer termo-$\lambda$ definido
        nos slides das aulas sobre o cálculo lambda.

        \medskip

		\begin{solution}

        \noindent\textbf{Observação geral:}
		Em todos os itens, $\Gamma$ segue a estrutura dos slides:
		\begin{align*}
        	\Gamma\equiv\lambda f\,x.\,(\text{predicado})\,(\text{caso base})\,(\text{chamada recursiva})
		\end{align*}
		
		Termos usados dos slides: $T$, $F$, $\lnot$, $Z$ (teste de zero), $P$ (antecessor), $S$ (sucessor),
        $+$, $\times$ e $G$ ($Gxy = T$ se $x\geq y$).

		\end{solution}

    \begin{Qlist}
        \item[(a)] Dado um natural $n$, $(Y\Gamma)n$ é verdadeiro, se $n$ é par, ou falso, se $n$ é ímpar.

        \begin{solution}
        $0$ é par, e $n>0$ é par se, e somente se, $n-1$ é ímpar:
        \[
            \Gamma \;\equiv\; \lambda fn.\,(Z\,n)\,T\,\bigl(\lnot\,(f\,(P\,n))\bigr).
        \]
        Exemplo: $(Y\Gamma)\,2 \equiv \lnot((Y\Gamma)\,1) \equiv \lnot(\lnot((Y\Gamma)\,0))
        \equiv \lnot(\lnot T) \equiv T$.
        \end{solution}

        \item[(b)] Dados dois naturais $a$ e $b$, $(Y\Gamma)ab$ é o resto da divisão euclidiana de $a$ por $b$.

        \begin{solution}
        Se $a\geq b$, o resto de $a$ por $b$ é igual ao resto de $a-b$ por $b$; caso contrário, é o próprio $a$:
        \[
            \Gamma \;\equiv\; \lambda fab.\,(G\,a\,b)\,\bigl(f\,(bPa)\,b\bigr)\,a.
        \]
        Exemplo: $(Y\Gamma)\,7\,3 \equiv (Y\Gamma)\,4\,3 \equiv (Y\Gamma)\,1\,3 \equiv 1$, pois $G\,1\,3 = F$.

        \end{solution}

        \item[(c)] Dados dois naturais $a$ e $b$, $(Y\Gamma)ab$ é o quociente da divisão euclidiana de $a$ por $b$.

        \begin{solution}
        Se $a\geq b$, o quociente é um a mais que o quociente de $a-b$ por $b$; caso contrário, é $0$:
        \[
            \Gamma \;\equiv\; \lambda fab.\,(G\,a\,b)\,\bigl(S\,(f\,(bPa)\,b)\bigr)\,0.
        \]
        Exemplo: $(Y\Gamma)\,7\,3 \equiv S((Y\Gamma)\,4\,3) \equiv S(S((Y\Gamma)\,1\,3))
        \equiv S(S\,0) \equiv 2$.

        \end{solution}

        \item[(d)] Dados dois naturais $a$ e $b$, $(Y\Gamma)ab$ é o maior divisor comum $(a, b$) entre $a$ e 
            $b$. Dica: $(a, b) = (b, a - b)$ e $(a, 0) = a$.

        \begin{solution}
        Caso base: $(a,0) = a$. Como a subtração $bPa$ resulta em $0$ quando $a<b$, a relação $(a,b)=(b,a-b)$ só pode ser aplicada diretamente se $a\geq b$; quando $a<b$, basta inverter a ordem dos argumentos devido a simetria $(a,b)=(b,a)$:
        \[
            \Gamma \;\equiv\; \lambda fab.\,(Z\,b)\,a\,\Bigl((G\,a\,b)\,\bigl(f\,b\,(bPa)\bigr)\,(f\,b\,a)\Bigr).
        \]
        Exemplo: $(Y\Gamma)\,4\,6 \equiv (Y\Gamma)\,6\,4 \equiv (Y\Gamma)\,4\,2 \equiv (Y\Gamma)\,2\,2
        \equiv (Y\Gamma)\,2\,0 \equiv 2$.
        \end{solution}

        \item[(e)] Dado um natural $n$, $(Y\Gamma)n = n!!$, onde 
        \[
            n!! = n(n - 2)(n - 4)\ldots .. 1
        \]
        e $1!! = 0!! = 1$.

        \begin{solution}
        Os casos base são $n=0$ e $n=1$, os quais satisfazem $Pn = 0$ (pois $P0=0$ e $P1=0$); para $n\geq 2$, temos $n!! = n\times(n-2)!!$:
        \[
            \Gamma \;\equiv\; \lambda fn.\,\bigl(Z\,(P\,n)\bigr)\,1\,\Bigl(\times\,n\,\bigl(f\,(P\,(P\,n))\bigr)\Bigr).
        \]
        Exemplo: $(Y\Gamma)\,5 \equiv \times\,5\,((Y\Gamma)\,3) \equiv \times\,5\,(\times\,3\,((Y\Gamma)\,1))
        \equiv \times\,5\,(\times\,3\,1) \equiv 15$.
        \end{solution}

        \item[(f)] Dados dois naturais $a$ e $b$, $(Y\Gamma)ab$ é igual a $a^b$.

        \begin{solution}
        Caso base: $a^0 = 1$. Para $b>0$, $a^b = a\times a^{b-1}$:
        \[
            \Gamma \;\equiv\; \lambda fab.\,(Z\,b)\,1\,\bigl(\times\,a\,(f\,a\,(P\,b))\bigr).
        \]
        Exemplo: $(Y\Gamma)\,2\,3 \equiv \times\,2\,(\times\,2\,(\times\,2\,1)) \equiv 8$.
        \end{solution}

        \item[(g)] Dado um natural $n$, $(Y\Gamma)n$ é a soma dos primeiros $n$ números naturais.

        \begin{solution}
        A soma de zero termos é $0$, e $\sum_{i=1}^{n} i = n + \sum_{i=1}^{n-1} i$:
        \[
            \Gamma \;\equiv\; \lambda fn.\,(Z\,n)\,0\,\bigl(+\,n\,(f\,(P\,n))\bigr).
        \]

		Considerando que o menor natural é $1$.

        Exemplo: $(Y\Gamma)\,3 \equiv +\,3\,(+\,2\,(+\,1\,0)) \equiv 6$.
        \end{solution}

        \pagebreak
        \item[(h)] Dado um natural $n$, $(Y\Gamma)n$ é a soma dos quadrados dos primeiros $n$ números naturais.

        \begin{solution}
        Análogo ao item anterior, somando $n^2 = \times\,n\,n$:
        \[
            \Gamma \;\equiv\; \lambda fn.\,(Z\,n)\,0\,\bigl(+\,(\times\,n\,n)\,(f\,(P\,n))\bigr).
        \]
        Exemplo: $(Y\Gamma)\,3 \equiv +\,9\,(+\,4\,(+\,1\,0)) \equiv 14$.
        \end{solution}

        \item[(i)] Dados três naturais $a, r$ e $n$, $(Y\Gamma)arn$ é igual a soma dos primeiros $n$ termos
            de uma progressão aritmética cujo primeiro termo é igual a $a$ e a razão é igual a $r$.

        \begin{solution}
        Os $n$ primeiros termos são $a, a+r, \ldots, a+(n-1)r$. A soma é $0$ se $n=0$; caso contrário, é o termo atual $a$ mais a soma dos $n-1$ primeiros termos da PA que se inicia em $a+r$ com a mesma razão:
        \[
            \Gamma \;\equiv\; \lambda farn.\,(Z\,n)\,0\,\bigl(+\,a\,(f\,(+\,a\,r)\,r\,(P\,n))\bigr).
        \]
        Exemplo ($a=1$, $r=2$, $n=3$): $(Y\Gamma)\,1\,2\,3 \equiv +\,1\,(+\,3\,(+\,5\,0)) \equiv 9$.
        \end{solution}

        \item[(j)] Dados três naturais $a, q$ e $n$, $(Y\Gamma)aqn$ é igual a soma dos primeiros $n$ termos
            de uma progressão geométrica cujo primeiro termo é igual a $a$ e a razão é igual a $q$.

        \begin{solution}
        Os $n$ primeiros termos são $a, aq, \ldots, aq^{n-1}$. De forma análoga à PA, a soma é $0$ se $n=0$; caso contrário, é $a$ mais a soma dos $n-1$ termos restantes da PG que começa em $a\times q$:
        \[
            \Gamma \;\equiv\; \lambda faqn.\,(Z\,n)\,0\,\bigl(+\,a\,(f\,(\times\,a\,q)\,q\,(P\,n))\bigr).
        \]
        Exemplo ($a=1$, $q=2$, $n=3$): $(Y\Gamma)\,1\,2\,3 \equiv +\,1\,(+\,2\,(+\,4\,0)) \equiv 7$.
        \end{solution}
    \end{Qlist}
} }

\end{document}

